% English translation of questions/5776/semB-moedC/q1.tex (metadata is taken from the Hebrew file).

\begin{question}
\textbf{Prove or disprove:} The following equation has a solution in $[0,\pi]$:
\[ \ln(x+1)-\sin(x)=1 \]
\end{question}

\begin{hint}
Define $g(x)=\ln(x+1)-\sin(x)-1$ and check its sign at the endpoints of the interval.
\end{hint}

\begin{solution}
\textbf{The statement is true.} Define $g(x)=\ln(x+1)-\sin x-1$. The function $g$ is continuous on $[0,\pi]$ as a sum of continuous functions ($\ln(x+1)$ is continuous for $x>-1$).

At the endpoints:
\[ g(0)=\ln1-\sin0-1=-1<0,\qquad g(\pi)=\ln(\pi+1)-\sin\pi-1=\ln(\pi+1)-1>0, \]
where the last inequality follows from $\pi+1>e$ (since $\pi+1>4>e$), and hence $\ln(\pi+1)>\ln e=1$. (Numerically $\ln(\pi+1)\approx1.421$.)

By the Intermediate Value Theorem (Bolzano's theorem) there exists $c\in(0,\pi)$ for which $g(c)=0$, i.e., $\ln(c+1)-\sin c=1$. Hence the equation has a solution in the interval.
\end{solution}
